% GENERATED FILE. Edit canonical lessons, then run npm run generate:books.
% canonical-source-hash: fnv1a64:883ef404ebfa7899
% canonical-lessons: RU-W81-yo, RU-R81-letters-first-pass, RU-R81-letters-second-pass

\chapter{The Letter Yo}
\label{ch:ru-letter-yo}
\hlchaptermodality{81}

\begin{chapteropening}
\textbf{By the end of this chapter:} \emph{I can write yo, the letter with two dots, and find it in words I already say.}

Complete the last lesson of chapter 81: I can write yo, the letter with two dots, and find it in words I already say.
\end{chapteropening}

\section[\texorpdfstring{\ru{ё} (yo)}{ё (yo)}]{\ru{ё} (yo) — the yo inside \ru{ребёнок}}

\label{lesson:RU-W81-yo}



\begin{quote}
Before the new one: write \ru{я} and \ru{ю} once each, and name them.

Now say \textbf{\ru{ребёнок}} (\emph{rebyónok}), \textquotedblleft{}a child\textquotedblright{}. You have read it for a long time. This lesson takes one letter out of it and puts it in your hand.
\end{quote}

\subsection*{Script you'll notice: \ru{ё}}
\textbf{\ru{ё}} — \emph{yo}.

It is inside \textbf{\ru{ребёнок}} (\emph{rebyónok}), \textquotedblleft{}a child\textquotedblright{}, and in \textbf{\ru{мёд}}, \textbf{\ru{чёрный}} and \textbf{\ru{тёплый}}. It is \textbf{\ru{е}} with \textbf{two dots} above it, and the syllable that carries it takes the stress. Many printed texts leave the dots off; this book keeps them, so you can see how the word is said.

\subsection*{Writing: \ru{ё}}
\hlblockfigure{\detokenize{figures/RU-W81-yo-filmstrip.pdf}}{How \ru{ё} is written, stroke by stroke}

Follow the numbered strip: start where movement 1 begins, and let each panel show you the next movement before your pen makes it. Then write \ru{ё} yourself — slowly, and larger than it is printed.

\begin{quote}
The strip shows \textbf{where to start this letter and which way to travel}, in one attested order, and names its source beneath it. Teachers vary, so treat it as a sound way in, not the only one.
\end{quote}

\subsection*{Your turn}
\begin{itemize}
\raggedright
  \item \emph{Look:} at \textbf{\ru{ребёнок}} and point at \ru{ё} inside it
  \item \emph{Trace it:} \ru{ё} three times, saying \emph{yo} as you finish each one
  \item \emph{Write it:} \ru{ю}, then \ru{ё}, from memory
  \item \emph{Read it:} \textbf{\ru{ребёнок}} aloud once more
\end{itemize}

\begin{tcolorbox}[breakable,colback=teal!4,colframe=teal!35!black,title={Before you move on}]
Which letter is this — \ru{ё}? (\textbf{yo}, \emph{yo}.) Which word did it come from? (\textbf{\ru{ребёнок}}, \emph{rebyónok}, \textquotedblleft{}a child\textquotedblright{}.)
\end{tcolorbox}
\section[(one letter)]{(one letter) — From the letter to the word — a first pass}

\label{lesson:RU-R81-letters-first-pass}



\begin{quote}
Say \textbf{\ru{ребёнок}} (\emph{rebyónok}), then write the letter it gave you.
\end{quote}

\subsection*{Your turn}
\noindent
\begin{tabularx}{\linewidth}{@{}>{\raggedright\arraybackslash}X>{\raggedright\arraybackslash}X>{\raggedright\arraybackslash}X@{}}
\toprule
\textbf{letter} & \textbf{name} & \textbf{the word it came from} \\
\midrule
\ru{ё} & \emph{yo} & \textbf{\ru{ребёнок}} \emph{rebyónok} \\
\bottomrule
\end{tabularx}

\begin{itemize}
\raggedright
  \item \emph{Write it:} each letter in the table from memory, then check it
  \item \emph{Say it:} the word each one came from
  \item \emph{Read it:} each word in the table aloud, and point at its letter
\end{itemize}

\begin{tcolorbox}[breakable,colback=teal!4,colframe=teal!35!black,title={Before you move on}]
Which word did each letter come from? (\textbf{\ru{ё}} from \textbf{\ru{ребёнок}}.)
\end{tcolorbox}
\section[(one letter)]{(one letter) — From the word back to the letter — a second pass}

\label{lesson:RU-R81-letters-second-pass}



\begin{quote}
Write \textbf{\ru{ё}} without looking, then say the word it came from.
\end{quote}

\subsection*{Your turn}
\noindent
\begin{tabularx}{\linewidth}{@{}>{\raggedright\arraybackslash}X>{\raggedright\arraybackslash}X>{\raggedright\arraybackslash}X@{}}
\toprule
\textbf{letter} & \textbf{name} & \textbf{the word it came from} \\
\midrule
\ru{ё} & \emph{yo} & \textbf{\ru{ребёнок}} \emph{rebyónok} \\
\bottomrule
\end{tabularx}

\begin{itemize}
\raggedright
  \item \emph{Say it:} each word in the table, then write the letter it gave you without looking
  \item \emph{Write it:} each letter once more, slowly, larger than it is printed
\end{itemize}

\begin{tcolorbox}[breakable,colback=teal!4,colframe=teal!35!black,title={Before you move on}]
Say each word in the table once more, and name the letter you took from it.
\end{tcolorbox}
